\documentclass[10pt,a4paper]{amsart}
\usepackage[margin=19mm]{geometry}
\usepackage{amsmath,amssymb,mathtools}
\usepackage[scr=rsfs,cal=euler]{mathalfa}
\usepackage{xfrac}
\usepackage[unicode,hidelinks,bookmarks=false]{hyperref}
\newcommand{\ie}{i.e.\ }
\newcommand{\cf}{cf.\ }
\newcommand{\resp}{respectively\ }
\newcommand{\Subsection}{\subsection}
\providecommand{\nobreakdash}{\nobreak\hbox{-}}
\setlength{\emergencystretch}{2em}
\multlinegap0pt
\usepackage{float}
\usepackage{tikz}
\usepackage{mathtools}
\usepackage{amssymb}
\usepackage{xcolor}
\usepackage{csquotes}
\newtheorem{remark}{Remark}
\newtheorem{proposition}{Proposition}
\title{On the Donaldson-Scaduto conjecture}
\author{Saman Habibi Esfahani, Yang Li, Department of Mathematics, Duke University, 120 Science Dr, Durham, NC 27708-0320, Department of Mathematics, Massachusetts Institute of Technology, 77 Massachusetts Avenue, Cambridge, MA 02139}
\date{Source: arXiv:2401.15432v2 (CC BY 4.0)}
\begin{document}
\maketitle

\noindent\emph{Source and licence.} On the Donaldson-Scaduto conjecture. Version arXiv:2401.15432v2 (CC BY 4.0), distributed by arXiv under the Creative Commons Attribution 4.0 International licence, http://creativecommons.org/licenses/by/4.0/, as recorded in the arXiv licence metadata for this exact version and shown on the abstract page https://arxiv.org/abs/2401.15432v2. Full licence notice: \texttt{licenses/arxiv-2401.15432-CC-BY-4.0.txt}.\par\smallskip

\noindent\textit{Editorial scope.} Counted unit: the article's proposition prop:gradientdivergence (source0.tex lines 457--465). The article's complete original proof of it is delivered with it (source0.tex lines 468--514, one \texttt{proof} environment of 47 lines). No mathematics is written, repaired or re-derived: the statement and the proof are the article's own lines, in the article's own notation, and no retained line is substituted or re-flowed.  Retained uncounted as the source-local context the proof needs: lines 334--338; lines 437--439.  One declared adaptation: exactly 1 ancillary strings inside the retained spans are omitted (image-inclusion commands and, where the source repeats a label, the repeated label definitions) and no other line differs from the source; the figure environments, with their placement, captions and labels, are kept, and the package records the omitted strings in fidelity receipts bound to the affected bodies.  No retained line contains a citation, so the wrapper carries no bibliography and no bibliography entry.  The retained lines contain the article's own diagram code, which the wrapper typesets with the standard tikz packages; no image file is delivered and the package declares no asset dependency.  Editorial formatting only, in addition: an AMS \texttt{amsart} preamble that restates the article's own declarations and macros used by the retained lines, namely the environments \texttt{remark}, \texttt{proposition} on separate counters, together with the standard packages those lines need.  The retained environments print in this document as Remark 1, Proposition 1 in order of appearance, whereas the article prints the same items with the numbering of its own full document.  The complete native source of the article as distributed in the arXiv e-print for this exact version is bundled unchanged as \texttt{sources/153607/source0.tex}.
\begin{remark}\label{standardposition}
The real Monge-Amp\`ere equation is invariant under adding an affine linear function. Geometrically, adding a constant to $\varphi$ has no effect on the special Lagrangian $L$, and adding a linear function amounts to translating $L$ along some vector in $\mathbb{R}_{(y_1,y_2)}^2$. In particular, for the triangle case $n=3$, we can reduce to the zero boundary data.  

The real Monge-Amp\`ere equation is also invariant under the Euclidean motion of the convex polygon $U$, with the corresponding change to $V$. Later, when we analyze the local behavior of $\varphi$ near an open edge, we will sometimes reduce to the `standard position' where the open edge is contained in the $u_2$-axis, and $U$ lies inside the right half-plane, to simplify the notations.
\end{remark}

\begin{equation}\label{boundarycontinuity}
\bar{\phi}(u)- C\text{dist}(u, \partial U)^{1/2} \leq \varphi(u)\leq \bar{\phi}(u)+ C\text{dist}(u, \partial U). 
\end{equation}

\begin{proposition}\label{prop:gradientdivergence}
\label{blowup}
Let $u_*$ be any point on a boundary edge $(p_i, p_{i+1})$ in $\partial U\setminus \{ p_1,\ldots, p_n\}$. Then, as $u\in U$ tends to $u_*$, the normal and tangential gradient components satisfy
\begin{alignat*}{3}
        \nabla^{\text{normal}} \varphi (u) &\to +\infty, 
        \quad
         \nabla \varphi (u)\cdot (p_{i+1}-p_i) &\to c_i.
\end{alignat*}
\end{proposition}

\begin{proof}\label{blowupproof}
The tangential component converges to a constant by the convexity of $\varphi$, the affine linearity of the boundary data, and the boundary continuity estimate (\ref{boundarycontinuity}) by considering the convex function restricted to line segments parallel to the boundary edge. 

We now focus on the normal gradient component and place the convex polygon into the standard position by Remark \ref{standardposition}, so the open boundary edge $(p_i, p_{i+1})$ containing $u_*$ lies on the $u_2$-axis, the domain is contained in the right half-plane, and the boundary data is zero on this edge. The normal gradient component is just $-\partial_{u_1}\varphi$. Using convexity, $\partial_{u_1}\varphi$ is bounded from the above near $u_*$.  








We suppose for contradiction that $\partial_{u_1} \varphi (u)$ stays bounded for some sequence $u \to u_*$, and therefore, the gradient at $u$ stays bounded. Then there exists some subgradient for $\varphi$ at $u_*$, which must be of the form $(-\Lambda, 0)$ for some $\Lambda\in \mathbb{R}$, because the tangential component is zero. 
Let 
\begin{align*}
    \Lambda_0 = \inf \{ \Lambda \; | \; (-\Lambda,0) \text{ is a subgradient of } \varphi \text{ at some point } v \in (p_i,p_{i+1})  \},
\end{align*}
so in particular $\varphi(u)\geq -\Lambda_0 u_1$ for any $u\in \overline{U}$. Thus $(-\Lambda_0, 0)$ is a subgradient at every point on the open edge, and $\partial_{u_1} \varphi\geq -\Lambda_0$ on $U$.




We fix a small constant $h > 0$ such that $u_*$ has distance at least $2h$ to the vertices. Let $R(u_*, h, \varepsilon)$ be the rectangle with length $h$ and width $\varepsilon\ll 1$, as shown in Figure \ref{rectangle}.
    \begin{figure}[H]
    
    \centering
    \caption{Standard position in the case $n=3$.\\ $R(u_*,h, \varepsilon)$ is the shaded region.}
    \label{rectangle}
    \end{figure}
By the Monge-Amp\`ere equation, and the strict positivity of $V$,
\begin{align}\label{equationV}
  C^{-1}h\varepsilon  \leq \int_{R(u_*,h, \varepsilon)} V du_1 du_2 = \int_{\nabla \varphi(R(u_*,h, \varepsilon))} dy_1 dy_2.
\end{align}
By the convexity of $\varphi$ and the bound $-\Lambda_0 u_1\leq \varphi\leq Cu_1$, we deduce
\begin{align*}
    |\partial_{u_2} \varphi(u) | \leq C (\Lambda_0,h) u_1, \quad \text{ for all } \quad u \in R(u_*,h, \varepsilon).
\end{align*}
Thus by considering the gradient image,
\begin{align*}
    |\int_{\nabla \varphi(R(u_*,h, \varepsilon))} dy_1 dy_2| \leq C( \Lambda_0, h )\varepsilon  \sup_{ R(u_*,h, \varepsilon) }(\partial_{u_1} \varphi(u) + \Lambda_0),
\end{align*}
Contrasting with (\ref{equationV}), for any small $\varepsilon>0$,
\begin{align*}
   \sup_{ R(u_*,h, \varepsilon) }(\partial_{u_1} \varphi(u) + \Lambda_0)  \geq C(\Lambda_0, h)^{-1}.
\end{align*}
In the limit $\varepsilon\to 0$, we can extract some subgradient at some boundary point, which contradicts the minimality of $\Lambda_0$.
\end{proof}

\end{document}
